% DRAFT (2026-10-22, to-do step 32). Not loaded by main.tex until the Globe Explorer is complete. Every number below is printed by the
% chain itself in a committed file: website/art/out/explorer/{healthy-blood,planted-blood,purified,noise}.json (the globes run_sample.py
% drew) and website/art/out/sphere/sphere.json (run_sample.py --build atlas-sphere).
% Still to write before this chapter goes into the book (each from the one command's own output):
%   - the information sphere (after its test on IAM-Atlas v2.1's healthy cells; DEV_SPHERE_01.md)
%   - a tumor against the same person's own normal tissue (the tumor-pair globe, redrawn by the one command)
%   - IAM-A on the whole genome (DEV_IAMA_GLOBE_01.md, run 6 through the one command)
%   - what the panel cannot explain (the Deconvolver's leftover, step 30)
%   - each layer's shuffled null
%   - figures: screenshots of the chain's own pages from release check V13 (GLOBE_SHOTS)
%   - the atlas globes and the noise globe redrawn by run_sample.py --build atlas-sphere / explorer-noise and checked identical to the
%     site's files, as the purified globe was (sphere.json still names its old builder, website/art/make_sphere_3d.py)

\chapter{The Globe Explorer}\label{ch:explorer}

\section{What it is}

The Globe Explorer is a turning globe in a web page. Every globe in it is drawn by the chain's one command,
\texttt{run\_sample.py}: a reference globe by its build mode, a sample's globe by the reading itself. The Explorer only shows them;
it reads, tares and draws nothing of its own. The same drawing goes into every report, so a researcher sees a sample exactly as the
chain measured it.

\section{The address layout}

Every site of IAM-Atlas v2.1 has one fixed place on the globe, set by where the site sits on the genome: the chromosomes in order,
wrapped onto the sphere by a frozen rule (the chain's sky order). The sphere is cut into 786{,}432 equal pieces (HEALPix,
$N_\mathrm{side}=256$, the pixelization the microwave-sky maps use), and 813{,}950 of the atlas's 813{,}951 sites have a place. The
layout is the same on every globe, for every cell and every sample, so any two globes can be laid over each other pixel for pixel.

Four things follow from putting the genome on a sphere instead of a flat strip:
\begin{itemize}
\item \textbf{Equal area.} Every piece of the sphere is the same size, so no part of the genome is stretched.
\item \textbf{No edges.} There is no seam where the map starts and ends.
\item \textbf{Clustered change shows as patches.} Neighbors on the genome are neighbors on the globe, so a loss running along a
stretch of DNA draws as a patch, and scattered wear draws as speckle.
\item \textbf{The microwave sky's tools apply unchanged.} Masks, power spectra and the other methods of Chapter~\ref{ch:skytools} were
built for exactly this kind of map.
\end{itemize}
This globe is a map of the genome's addresses. The map of how the sites behave, the information sphere
(Section~\ref{sec:ag_infosphere}), is a different globe.

\section{The globes}

\paragraph{The atlas.} One globe for each of IAM-Atlas v2.1's 79 cell types: the atlas's best value at every site, and how sure
the atlas is of it.

\paragraph{Met-A's purified reference.} The reference Met-A and the C-Score read against: purified healthy arrays of eight blood
groups, at Met-A's 6{,}000 neutrophil identity sites. Between 5{,}548 and 5{,}553 of the 6{,}000 have a place on the layout, depending
on the group, so this globe is drawn as dots, each at the same place as on the atlas globes.

\paragraph{How sure the atlas is.} IAM-Atlas v2.1 was fitted the way a sky map is: every site carries 100 equally likely versions of
each cell's healthy map (posterior draws, deposited as Zenodo record 10.5281/zenodo.23422688). Averaging more draws beats the noise
down, about as one over the square root of their number, as more observing beats down the noise of the microwave sky. Sliding from 1
draw to 100 shows where the atlas is sure (it goes dark first) and where it is least sure (it stays lit). For neutrophils the median
spread is 0.0029 on the autosomes and 0.0105 on the X and Y; for monocytes, 0.0033 and 0.0103. The X and Y stay lit longest because
the atlas's donors are women and men.

\paragraph{A sample against the atlas.} A real healthy whole blood (GEO sample GSM7508429), read by the one command. Its healthy
reference is healthy blood mixed in this person's own cell shares, as the IAM-Deconvolver read them (47\,\% neutrophils, 23\,\% CD4 T
cells, 10\,\% CD8 T cells, 8\,\% NK cells, 6\,\% monocytes, 6\,\% B cells), compared at the 773{,}556 atlas sites the array measured.
The globe holds the sample, its healthy reference, the change between them, the change at Met-A's 6{,}000 sites against its purified
reference, and each cell's healthy map.

\paragraph{A constructed loss.} The same blood with a clustered loss planted at Met-A's neutrophil sites, labeled on its page as
constructed, not a person.

\section{The change, fixed by the measurement}

What counts as changed is decided by the measurement, never by the reader. Each site's difference from the healthy reference is
measured in that site's healthy width (how much healthy people vary there, and how sure the atlas is) and in the sample's own spread.
The differences are then read in blocks of 10 neighboring sites, as the C-Score reads them, and a block has changed when it passes the
line that pure chance crosses less than once in 20 genomes.

Blocks, because ordinary wear is scattered: it seldom moves a whole block, so a healthy sample stays nearly dark. A change that runs
along the genome moves whole blocks and lights up.

\begin{itemize}
\item The healthy blood: 45 of 77{,}355 blocks changed across the atlas, 1 of them on the X and Y; at Met-A's sites, 0 of 597.
\item The constructed loss: at Met-A's sites, 44 of 597 blocks changed. Across the whole atlas the count stays at 45 of 77{,}355,
because the loss was planted only at Met-A's 6{,}000 sites, which are spread thinly through the atlas's 773{,}556.
\end{itemize}
Each count is printed beside its layer and never moves.

\section{The slider}

The slider is a viewing tool. On a change layer it fades the unchanged background, so the change can be seen by itself; the changed
blocks and their count never move. On a map layer (a cell's methylation, or how sure the atlas is) it keeps only the sites at or above
the chosen value, and every other part of the globe goes dark: the lower sites, the sky with no site, and the rim of each dot alike.
At the right end of the slider, a map layer is dark.

\section{Checked in a browser}

The chain's release check opens every globe page in a browser, moves every slider on every layer as a reader would, and photographs
the globe at the left, middle and right of the slider (check V13). It requires that every layer has a slider; that every map layer is
at least 90\,\% dark at its slider's right end; that a healthy blood's change layers are dark at their right end; that the Met-A
layer's key says it shows disorder; and that the constructed loss lights the Met-A layer more than the healthy blood does.
